\chapter{Discusión y conclusiones}
\label{cap:discusion}

\section{Discusión de los resultados}

\subsection{Seguimiento de precios y capacidad de arbitraje (OE2, H1)}

El resultado de H1 es el más sólido del trabajo: el arbitraje activo reduce el área de exceso de error en todas las semillas y en todas las celdas, y la capacidad gradúa la recuperación, de $4.6\,\%$ sin arbitraje a $56.8\,\%$ con capacidad baja y a $100\,\%$ con capacidad alta. Es un resultado esperable desde la mecánica del modelo, porque el arbitrajista es el único agente que mueve el precio del AMM hacia la referencia, y por eso la contribución no está en el signo del efecto sino en su graduación. La tasa de recuperación con capacidad baja depende de la combinación entre el tamaño del choque y la profundidad (tabla~\ref{tab:recuperacion}): un choque grande requiere varios pasos acotados, y el horizonte de 20 ticks puede no bastar. Esa dependencia sugiere que, en este modelo, la capacidad de arbitraje relevante no es un umbral sino una razón entre el tamaño de la discrepancia y la capacidad por tick.

La lección de diseño es la precisión. La semiamplitud objetivo de H1 no se alcanzó porque el piloto combinó choques positivos y negativos cuyos efectos se cancelan casi por completo dentro de cada semilla (correlación $-0.981$). Un piloto construido con el mismo diseño que la confirmación habría anticipado una desviación nueve veces mayor. La conclusión de que el arbitraje reduce el exceso no depende de ello, pero la estimación del tamaño del efecto es menos precisa de lo previsto.

\subsection{Liquidez adaptativa (OE3, H2)}

H2 es positivo en el diseño confirmatorio, pero su interpretación de ingeniería está condicionada. En el análisis suplementario original de 40\,500 trayectorias, la diferencia de respuesta máxima entre E3 adaptativo y estático con profundidad 0.5 fue 0.03458 unidades logarítmicas con órdenes absolutas fijas (intervalo descriptivo del 95\,\% sin ajuste [0.03225, 0.03690]); con profundidades 1 y 2 fue efectivamente cero. Al escalar las órdenes de ruido y tendencia con las reservas iniciales, fue efectivamente cero en todas las profundidades. Desactivar la orden de pánico y mantener la tendencia redujo el contraste de poca profundidad a 0.01471 [0.01364, 0.01577]. La sensibilidad atribuye así el patrón original a la interacción entre el retiro, el tamaño relativo de la orden y la conducta de pánico, no a la profundidad como factor de riesgo operativo aislado. La ablación es posterior al protocolo y sus intervalos son descriptivos.

No se afirma que los proveedores reales amplifiquen las perturbaciones. La regla de retiro es una decisión de modelado, sin calibración y sin respaldo empírico propio. Su respaldo teórico es indirecto: la literatura trata la salida de los proveedores como un problema de decisión \cite{bergault2025exit} y documenta su comportamiento en intercambios descentralizados \cite{lehar2025uniswap}, pero ninguno de esos trabajos valida la regla usada aquí.

\subsection{Mapa dinámico reducido (OE4, H3)}

El resultado de H3 es mixto. El mapa de dos variables predice mejor que la persistencia, y con una mejora sustancial en las condiciones con arbitraje activo (15 a 21 por ciento en E1-alta y E3), pero no supera a un autorregresivo escalar y es ligeramente peor que la persistencia sin arbitraje. Dos observaciones limitan su interpretación. Primera, el segundo estado, la presión neta de compra, no mejora la predicción a un paso respecto del exceso de error solo; la ventaja del mapa sobre la persistencia proviene del coeficiente de contracción $a$, que un modelo escalar captura igual. Segunda, el coeficiente $\beta$ de E3 tiene signo opuesto al de las demás condiciones y la ecuación del exceso es degenerada en E1-máx, lo que muestra que la estructura de dos variables no es estable entre condiciones.

La clasificación de regímenes con el criterio de Schur no discrimina: todas las matrices ajustadas son estables, incluida la de E0, cuyo radio espectral es $0.998$ y cuyas ejecuciones son persistentes en $57\,\%$ de los casos. Esto es coherente con lo que el criterio puede decir (la estabilidad de una dinámica lineal autónoma en un punto) y no con el horizonte finito y el ruido del simulador. El mapa puede describir la contracción del error en algunas condiciones ajustadas, pero los resultados actuales no respaldan su uso como clasificador de regímenes ni como regla de control: no supera al autorregresivo escalar y sus etiquetas reservadas son débiles.

La comprobación cuadrática de Lyapunov precisa el alcance de este resultado. Para los mapas afines específicos de cada condición, una matriz $P$ definida positiva produce una disminución estricta de $V(d,z)$ cuando se fijan en cero la entrada y los residuos del ajuste. El límite de E0 está cerca de uno, mientras que los ajustes con arbitraje alto contraen más rápido en esta norma cuadrática. Esta es una propiedad matemática de los mapas autónomos ajustados, no un nuevo resultado empírico sobre el simulador. Tampoco se afirma que el estado registrado $S_t=(x_t,y_t,p_t^C,W_t)$ sea completo de Markov: la regla de tendencia usa el precio previo del AMM, así que una representación completa de transición debe conservar esa memoria, las banderas de regla no codificadas en otra parte y la entrada aleatoria o el estado de reproducción que corresponda.

Una comparación independiente de mapas sensibles a la banda produjo una reducción pequeña del MAE recursivo reservado del error de precio (aproximadamente 1.7\,\%), concentrada casi por completo en la condición de capacidad baja con profundidad 2. Los errores de presión de flujo no mejoraron. Por tanto, el candidato podría servir como modelo sustituto para pruebas de estrés específicas de una condición, pero no como modelo general de monitoreo. Un estudio empírico futuro podría evaluar si los registros sincronizados de la discrepancia de precio con signo respecto de una banda de costos, la capacidad o los límites de ejecución del arbitraje, la profundidad del pool y los retiros de liquidez mejoran la explicación o predicción. También podría comparar modelos sustitutos candidatos con mapas afines y referencias simples en datos cronológicos y separados por pool que incluyan comisiones, demoras de ejecución e inventarios. Estas son preguntas de validación propuestas; los resultados sintéticos no establecen requisitos de producción ni justifican el despliegue o el uso de umbrales operativos.

\section{Conclusiones por objetivo}

\begin{description}
  \item[OE1.] El simulador discreto superó nueve comprobaciones de implementación. Se reprodujeron con artefactos de resultados idénticos la confirmación original, la mecánica determinista, el barrido, la selección estocástica y el mapa reservado. El análisis de ingeniería separado usó 300 semillas nuevas y sus tablas se reprodujeron en otra carpeta de salida. La nueva comparación híbrida usó 48\,000 transiciones de selección y 9\,600 trayectorias reservadas con 300 semillas intactas. Estas comprobaciones verifican el código y la reproducibilidad, no la validez externa.
  \item[OE2.] El arbitraje activo reduce el área de exceso de error de $2.36$ a $0.10$, y la capacidad gradúa la tasa de recuperación en las condiciones simuladas.
  \item[OE3.] El retiro adaptativo aumenta la respuesta máxima en $0.011$ en promedio, pero solo con la menor profundidad ensayada, y aproximadamente la mitad de ese efecto se atribuye al agente de pánico. No se concluye amplificación con profundidades mayores.
  \item[OE4.] El mapa reducido supera en promedio a la persistencia, pero no al autorregresivo escalar; las etiquetas de régimen reservadas son débiles. Un mapa híbrido sensible a la banda reduce el MAE recursivo del error de precio en cerca de 1.7\,\% en una nueva cuadrícula sintética, pero casi toda la ganancia se concentra en arbitraje de capacidad baja y profundidad 2, mientras que no mejoran los errores de presión de flujo. El análisis no exhaustivo de Lyapunov certifica estabilidad asintótica solo para los mapas afines autónomos y sin entrada; no certifica el mapa híbrido ni el simulador completo. Estos resultados motivan registrar eventos y hacer pruebas de estrés específicas, no un clasificador o regla de control para producción.
\end{description}

\begin{samepage}
\section{Conclusión general}

En el micromundo ensayado, la capacidad de arbitraje gradúa la corrección de la discrepancia entre el AMM y la referencia, mientras que la profundidad cambia el impacto de una orden en relación con las reservas de la piscina. La sensibilidad con semillas nuevas añade una precisión importante: la amplificación aparente de poca profundidad desaparece cuando el tamaño de las órdenes escala con las reservas y disminuye al eliminar la respuesta de pánico. La recuperación también requiere medir persistencia: en la condición adaptativa original, el 97.4\,\% de las trayectorias volvió a la banda al menos una vez, el 81.9\,\% terminó dentro y el 31.2\,\% de las que se recuperaron salió después. El mapa reducido no superó al autorregresivo escalar y, con esta evidencia, no debe usarse para clasificar regímenes.

La implicación de ingeniería consiste en preguntas para futuros estudios empíricos, no en requisitos operativos validados. Un estudio CEX/DEX con datos reales podría evaluar si alinear temporalmente la cotización externa y el estado de la piscina, registrar reservas, tamaños de operación, cambios de liquidez, comisiones y demoras, expresar el tamaño de las órdenes en relación con las reservas disponibles y distinguir el primer retorno del seguimiento terminal y de las salidas posteriores mejora la explicación o predicción. También podría comparar modelos reducidos con la persistencia y un autorregresivo escalar en periodos cronológicos y piscinas separadas. Estas mediciones y comparaciones candidatas se proponen a partir de experimentos sintéticos; su valor operativo aún no se ha evaluado. El estudio confirmatorio del simulador no tiene entradas de mercado en vivo, comisiones, calibración ni validación externa. Un piloto separado registra solo tres intercambios DEX públicos y no aporta una comparación CEX/DEX sincronizada que pueda publicarse. Ninguno de los dos establece procedimientos ni umbrales de producción.

El piloto de datos añade una observación de ingeniería limitada: es posible ingerir eventos públicos del DEX conservando marcas de tiempo, orden, reservas, nivel de comisión del pool y estimaciones de deslizamiento provistas por la fuente. Su muestra de seis minutos y tres eventos comprueba la tubería, no constituye un hallazgo de mercado. Las observaciones derivadas del CEX se mantienen privadas mientras no se resuelvan sus derechos de redistribución, y el extracto público del DEX por sí solo no permite estimar el error de seguimiento entre mercados ni los costos reales de ejecución.

El resultado del mapa híbrido sugiere preguntas empíricas adicionales: si el arbitraje está limitado por capacidad, si se retiró liquidez y si los mapas sensibles a la banda difieren de los mapas afines entre regímenes operativos observados. La ganancia actual se concentra en un solo régimen sintético; aporta una hipótesis de diseño del modelo para futuros estudios de reproducción, no evidencia para desplegar Hyb. El error recursivo informado también usa como entrada la trayectoria realizada de choques futuros del CEX en el simulador. Una evaluación futura de pronóstico en tiempo real tendría que limitar el modelo a la información disponible en cada decisión o evaluar por separado el pronóstico del precio externo.

\end{samepage}

\section{Limitaciones y amenazas a la validez}

\begin{itemize}
  \item \textbf{Referencia exógena.} El CEX no responde al AMM, ni tiene profundidad, latencia o comisiones. Los efectos de realimentación entre mercados quedan fuera.
  \item \textbf{Sin comisiones.} El simulador solo implementa intercambios sin comisión. La banda de no arbitraje es un parámetro y no una consecuencia de un costo. Los niveles de comisión, de latencia y de volatilidad externa son candidatos para trabajos posteriores.
  \item \textbf{Calendario.} El orden de las acciones dentro de un tick afecta los resultados (por ejemplo, el agente de pánico observa el precio antes de que actúen el proveedor y el arbitrajista). No se evaluaron calendarios alternativos.
  \item \textbf{Definición de profundidad.} La profundidad escala las reservas y mantiene constante el tamaño absoluto de las órdenes. Otra definición, con órdenes proporcionales a la reserva, podría cambiar la conclusión sobre H2.
  \item \textbf{Inventarios.} Los inventarios son amplios y ninguna orden estuvo limitada por capital; las restricciones de capital reales no se representan.
  \item \textbf{Agentes de un solo tipo.} Cada clase de agente es representativa y no hay heterogeneidad dentro de una clase.
  \item \textbf{Solo disponibilidad de datos externos.} Las pruebas de hipótesis y la comparación de modelos usan datos sintéticos. Un piloto DEX separado, de seis minutos, contiene tres intercambios públicos y es demasiado pequeño y selectivo para evaluar el comportamiento del mercado; el manuscrito público no incluye valores derivados del CEX. No hay validación externa y no se evaluó la coherencia entre el simulador y los hechos estilizados de mercados reales.
  \item \textbf{Comprobación exploratoria.} La repetición de H2 sin agente de pánico no estaba preregistrada.
  \item \textbf{Generalización de H3.} La evaluación fuera de muestra cubre celdas y semillas no vistas, pero dentro del mismo simulador y de la misma familia de choques.
  \item \textbf{Alcance de estabilidad.} El certificado de Lyapunov se aplica a los mapas afines ajustados por condición, omitiendo la entrada y los residuos del ajuste. No se ha demostrado que el estado registrado sea completo de Markov ni se establece una función de Lyapunov o estabilidad estocástica para el simulador completo.
\end{itemize}
